\section{Acts over the monotone transformation monoids}
\label{sec:E}

Throughout this section $n\geq1$, $U_n=\Delta_{\leq n}$ and
$M_n=\End_\Delta([n])$. Multiplication is composition, $gf=g\circ f$;
a right act satisfies $(h\cdot g)\cdot f=h\cdot(gf)$. An acyclic
space is nonempty and has reduced integral homology zero in every degree.

\begin{claim}{E.1}{Proved}{All acts are uniquely represented by skeletal simplicial sets}
\label{E:acts}
Restriction to $[n]$ is an equivalence between presheaves on $U_n$ and
right $M_n$-acts. Consequently every small right act $H$ is isomorphic
to $K_n$ for a simplicial set $K$ with no nondegenerate simplices above
dimension $n$, unique up to isomorphism. The act is finite exactly when
$K$ has finitely many nondegenerate simplices.
\end{claim}
\begin{proof}
Choose monotone maps $i_k:[k]\to[n]$, $i_k(j)=j$, and
$r_k:[n]\to[k]$, $r_k(j)=\min(j,k)$, for $0\leq k\leq n$.
Then $r_ki_k=\id$ and $e_k=i_kr_k$ is an idempotent of $M_n$.
Given $H$, put $K_k=\{h\in H:h\cdot e_k=h\}$. For
$\alpha:[k]\to[l]$ define the set map
\[
 K(\alpha):K_l\longrightarrow K_k,
 \qquad h\longmapsto h\cdot(i_l\alpha r_k).
\]
The image is fixed by $e_k$. For $\beta:[l]\to[m]$, the identity
$r_li_l=\id$ gives
$K(\alpha)K(\beta)=K(\beta\alpha)$. The map assigned to an identity
is the identity on its fixed-point set. Also $K_n=H$ since $e_n=\id$.
Conversely, for a presheaf $K$ the split maps $K(r_k)$ and $K(i_k)$
identify $K_k$ with this fixed-point subset of $K_n$, naturally in $k$.
The same formulas extend morphisms of acts and prove the equivalence.
They give explicitly the usual splitting of idempotents; we do not
assert that the one-object category itself is equivalent to $U_n$.

The inclusion $U_n\hookrightarrow\Delta$ has restriction as a functor
on presheaves. Its left adjoint is skeletal extension. It is fully
faithful, and its image consists of the simplicial sets with all
nondegenerate simplices in degrees at most $n$. Indeed the
Eilenberg--Zilber decomposition writes each simplex uniquely as a
degeneracy of a nondegenerate simplex. In degrees at most $n$ this
decomposition follows from the simplicial identities already in
$U_n$. In higher degrees skeletal extension adjoins precisely those
degeneracies, with the face maps determined by the given truncated
data. This proves existence and uniqueness. Each $K_k$, $k\leq n$,
is a subset of $H$ after the displayed splitting; hence finite $H$
gives finitely many nondegenerate simplices. Conversely a finite
list of nondegenerate simplices has only finitely many degeneracies
in degree $n$. This includes the empty act and empty simplicial set.
\end{proof}

\begin{claim}{E.2}{Proved}{The comparison which holds for every simplicial set in scope}
\label{E:comparison}
For every simplicial set $K$ of dimension at most $n$, the inclusion
$j:M_n/K_n\to U_n/K$ is homotopy final and is a weak equivalence.
There is a natural zigzag of weak equivalences
\[
 B(M_n/K_n)\xrightarrow{Bj}B(U_n/K)
 \longleftarrow T_n(K)\longrightarrow |K|.
\]
The statement has no finiteness or nondegenerate-face hypothesis.
\end{claim}
\begin{proof}
An object $c=(k,x)$ of $U_n/K$ is a retract of
$d=(n,K(r_k)x)$, with arrows $i:c\to d$, $r:d\to c$ induced
by $i_k,r_k$ and $ri=\id_c$. In $c\backslash j$ the chosen object
is $(d,i)$. For an arbitrary object $(d',f:c\to d')$, the arrow
$fr:d\to d'$ is a comma morphism from $(d,i)$ to $(d',f)$.
These morphisms are natural in $(d',f)$. Thus the constant functor
at $(d,i)$ has a natural transformation to the identity. Each comma
category is contractible, proving the assertion by Theorem A.
This is exactly the retract argument of PRETOT A.4, and it uses no
assumption about faces of nondegenerate simplices.

Define $T_n(K)$ to be the homotopy colimit of the standard geometric
simplices over $U_n/K$: on $(k,x)$ its value is $|\Delta[k]|$, and
on a simplex map $\alpha$ its value is $|\alpha|$. Projection to
the classifying space of the indexing category is a weak equivalence,
since all these standard simplices are contractible. Evaluation of
each labelled simplex in $|K|$ gives the second map.

We verify that this evaluation is a weak equivalence. Both the bar
model of $T_n$ and realization preserve coproducts and pushouts of
simplicial sets, and send a monomorphism to a cofibration in their
realization models. More explicitly, degree $p$ of the bar model is
the coproduct, over chains $[k_0]\to\cdots\to[k_p]$ in $U_n$, of
$K_{k_p}\times|\Delta[k_0]|$. Thus it preserves those pushouts
degreewise; its degeneracy inclusions are inclusions of coproduct
summands. Pushouts along monomorphisms are consequently homotopy
pushouts after realization. For $K=\Delta[m]$, $m\leq n$, the
category $U_n/\Delta[m]$ has terminal object $([m],\id)$, so
both sides are contractible and evaluation is a weak equivalence.
Induct over the skeletal attachments
$\partial\Delta[m]\hookrightarrow\Delta[m]$, $0\leq m\leq n$.
The gluing property just proved establishes the result on each
skeleton, with arbitrary coproducts of cells. This proves the
assertion for infinite $K$ as well. For the empty $K$ every space
and category displayed is empty and the comparison remains valid.
\end{proof}

\begin{claim}{E.3}{Refuted}{The nondegenerate-simplex arrow in Lead 4(a)}
\label{E:ndcounter}
For $n\geq2$, the second proposed functor
$U_n/K\to\Delta_{\rm inj}/K^{\rm nd}$ from PRETOT A.4 need not
exist with the stated formula and its proposed target need not have
the homotopy type of $K$. Take the finite simplicial set
$K=\Delta[2]/\partial\Delta[2]$.
For $n=1$ the assertion holds, as it also does for any $n$ when $K$
is a free degeneracy completion of a semisimplicial set.
\end{claim}
\begin{proof}
The nondegenerate simplices of this $K$ are one vertex $v$ and one
two-simplex $t$. Each face of $t$ is the degenerate edge $s_0v$.
Thus nondegenerate simplices do not form a semisimplicial set. Even
if the proposed notation is repaired to mean the category whose
objects are the nondegenerate simplices and whose arrows are
injections satisfying the simplex equation, that category has two
objects, three arrows $v\to t$, and their identities. Its nerve
is a graph with two vertices and three edges, so $H_1=\Z^2$ and
$H_2=0$. In contrast, $|K|=S^2$.

The precise functoriality failure occurs on a face arrow
$(1,s_0v)\to(2,t)$. The factorization asserted in PRETOT A.4
would express its injective map $[1]\to[2]$ as a map through the
degeneracy quotient $[1]\twoheadrightarrow[0]$, which is impossible.
The earlier return's actual example is a free degeneracy completion,
where every face is still a specified nondegenerate simplex before
degeneracies are added; its argument is valid there. In dimension
one all faces of nondegenerate edges are vertices, so that same
argument applies to every simplicial set of dimension at most one.
The program \texttt{code/E\_acts\_checks.py} verifies the finite
quotient, monoid action, category counts and integer boundary ranks.
Claim E.2 supplies the correct comparison for all dimensions.
\end{proof}

\begin{claim}{E.4}{Proved}{Classification of totally aspherical acts}
\label{E:total}
Let $H=K_n$ under E.1. The following are equivalent:
\begin{enumerate}[label=(\roman*)]
\item $H$ is a totally $W_{\Z}$-aspherical right $M_n$-act;
\item $|K|$ is acyclic;
\item $K$ is nonempty and connected and the right $\Z M_n$-module
      $\Z[K_n]$ is projective.
\end{enumerate}
Moreover $H$ is totally $W_\infty$-aspherical if and only if $|K|$
is weakly contractible. These equivalences hold for finite and
infinite acts.
\end{claim}
\begin{proof}
The identity test in totality and E.2 prove (i)$\Rightarrow$(ii).
For (ii)$\Rightarrow$(iii), normalize the simplicial abelian group
$\Z[K]$. Its chain groups $C_k$ are free on the nondegenerate
$k$-simplices and vanish for $k>n$. Write $B_k=\im(d:C_{k+1}\to C_k)$.
Acyclicity and freeness of subgroups of free abelian groups give
split exact sequences
\[
 0\to B_k\to C_k\to B_{k-1}\to0\quad(k>0),\qquad
 0\to B_0\to C_0\to\Z\to0.
\]
Consequently there is an actual chain isomorphism
\[
 C\cong S^0(\Z)\oplus\bigoplus_{k=1}^n D^k(B_{k-1}),
\]
where $D^k(P)$ has differential $\id_P$ in degrees $k,k-1$.
Both $S^0(\Z)$ and $D^k(P)$ for free $P$ are projective in
the abelian category of complexes concentrated in $[0,n]$:
maps from them are respectively evaluation in degree zero or
degree $k$, followed by a free-group Hom functor. The Dold--Kan
equivalence, restricted to degrees at most $n$, is an equivalence
of abelian categories and preserves projectives. This restriction
follows directly from its formula
$\Gamma(C)_m=\bigoplus_{[m]\twoheadrightarrow[k]}C_k$;
see \citet[Theorem 14.24.3]{StacksDK}. Therefore
$\Z[K]|_{U_n}$ is projective. The idempotent-splitting equivalence
in E.1 works verbatim for abelian-group-valued presheaves, and
preserves projectives. It gives (iii).

For completeness we verify the precise use of PRETOT A.1.
For a covariant act $L:M_n\to\Set$, the two-sided integral bar
complex for $\Z[H]$ and $\Z[L]$ has a preferred basis equal to the
simplices of the restricted action category
$\El_{M_n/H}(L\pi)$. A chain of base arrows, its final $H$-element,
and its initial $L$-element determine all intermediate elements.
The faces and degeneracies agree, so its homology is
$\Tor_i^{\Z M_n}(\Z[H],\Z[L])$. Projectivity makes this zero for
$i>0$, and each component is acyclic. For any functor $p:M_n\to U'$,
the canonical comparison
$M_n/H\to U'/p_!H$ is final. Its under-comma at $(u',c)$ is the
component $c$ of the restricted action for
$L(m)=U'(u',p(m))$. Homological Theorem A therefore makes this
comparison an integral homology equivalence. Its source is acyclic
by E.2 and (ii), or by the same bar calculation with $L=1$ and the
connectedness hypothesis in (iii). Thus (iii)$\Rightarrow$(i).

For weak-homotopy totality necessity again follows from the identity
test. Conversely assume $|K|$ weakly contractible. We have just
proved that every target $U'/p_!H$ is acyclic. A final functor
surjects on fundamental groups: one lifts each finite target loop
using chosen objects of its nonempty under-comma categories, and
connects the choices by zigzags in those connected comma categories.
The commuting triangles give the corresponding lifted loop.
Here the source is simply connected by E.2, so each target is
simply connected. The homological Whitehead theorem makes it weakly
contractible. This proves the assertion for every $p$, with no
replacement of restricted actions by arbitrary actions on $M_n/H$.
\end{proof}

\begin{claim}{E.5}{Proved}{Fundamental groups for finite and arbitrary acts}
\label{E:groups}
The groups occurring as $\pi_1B(M_n/H)$ for finite totally
$W_{\Z}$-aspherical right acts are precisely
\[
\begin{array}{c|l}
n=1&\{1\},\\
n=2&\text{groups admitting a finite balanced presentation with trivial abelianization},\\
n\geq3&\text{finitely presented groups }G\text{ with }H_1(G;\Z)=H_2(G;\Z)=0.
\end{array}
\]
The middle class is equivalently the fundamental groups of finite
acyclic two-dimensional CW complexes. Without finiteness, the
answer is respectively the trivial group; the groups admitting a
presentation whose exponent-sum homomorphism
$\Z^{(R)}\to\Z^{(X)}$ is an isomorphism; and all small groups
$G$ with $H_1(G;\Z)=H_2(G;\Z)=0$.
\end{claim}
\begin{proof}
By E.1, E.2 and E.4 it suffices to classify fundamental groups
of acyclic simplicial sets of dimension at most $n$, with finitely
many nondegenerate simplices when finiteness is required. Such a
realization is a CW complex with exactly those cells. An acyclic
connected graph is a tree up to homotopy, so its fundamental group
is trivial, also for infinite graphs.

For a connected acyclic two-complex choose a maximal tree in its
one-skeleton. Collapsing that tree gives a homotopy equivalent
presentation complex with generators $X$ and relators $R$.
Its cellular boundary $C_2\to C_1$ is precisely the exponent-sum
map $\Z^{(R)}\to\Z^{(X)}$. Acyclicity says exactly that this is
an isomorphism. Conversely a presentation with that property has
an acyclic presentation complex. For finite $X,R$, this is
equivalent to $|X|=|R|$ and trivial abelianization: a surjective
square integer matrix is unimodular. For infinite sets equality
of the two cardinalities alone is insufficient; the isomorphism
condition is retained as stated.

For any acyclic connected complex $X$, abelianization gives
$H_1(\pi_1X;\Z)=0$ and the Hopf exact sequence
$\pi_2X\to H_2(X;\Z)\to H_2(\pi_1X;\Z)\to0$ gives the
second vanishing. If $X$ is finite, its fundamental group is
finitely presented. Conversely take a presentation complex $Y$
of a superperfect group $G$. The group $H_2(Y;\Z)$ is free,
being a subgroup of the free group $C_2(Y;\Z)$. The Hopf sequence
says that every basis element is represented by a sphere
$S^2\to Y$. Attach one three-cell along each such representative.
The resulting complex $X$ has the same fundamental group, and
its cellular differential $C_3\to C_2$ is injective with image
$\ker(d:C_2\to C_1)$. Hence $H_i(X;\Z)=0$ for every $i>0$.
When the presentation is finite, the basis and the attachment
are finite. No finiteness is needed in the general construction.

We justify the passage to simplicial sets without assuming that
an arbitrary attaching map respects a global vertex ordering.
Here is a dimension-preserving mapping-cylinder construction.
For a simplicial map $f:L\to T$ of abstract simplicial complexes,
let $P(L),P(T)$ be their nonempty face posets. Form the poset
$P(f)=P(T)\amalg P(L)$ by retaining both face orders and adding
$\tau\leq\sigma$ for $\tau\in P(T)$, $\sigma\in P(L)$ precisely
when $\tau\subseteq f(\sigma)$. The retraction
$r:P(f)\to P(T)$ sends $\sigma$ to $f(\sigma)$ and is right
adjoint to the inclusion. Thus $Br$ is a homotopy equivalence,
and its restriction to $P(L)$ realizes the subdivision of $f$.
A mixed chain whose largest target face has dimension $r_0$ and
smallest domain face has dimension $q_0$ has $r_0\leq q_0$.
Counting its target and domain portions gives
\[
 \dim NP(f)\leq\max\{\dim T,\dim L+1\}.
\]
The inclusion $NP(L)\hookrightarrow NP(f)$ is a simplicial
subcomplex. Attaching its simplicial cone therefore models the
homotopy cofiber of $f$, in this same dimension bound.

Start with a simplicial graph of subdivided loops for the
generators. Simplicial approximation on each finite polygonal
attaching circle and the preceding construction realize the
presentation complex by a two-dimensional simplicial complex.
Use separately tagged domain face posets sharing $P(T)$ for
multiple attachments; no mixed chain passes from one domain
copy to another, so the dimension bound is unchanged. Next
represent the chosen elements of $\pi_2$ by maps of finite
triangulated two-spheres. Each compact image lies in a finite
subcomplex; simplicial approximation after domain subdivision
gives the required simplicial maps into the existing complex.
The same construction adds the three-cells without increasing
dimension beyond three. The nerve of every face poset is a
simplicial set, so the desired objects are obtained. Finiteness
is preserved when there are finitely many generators, relators,
and three-cells; the construction also works for any small
family of attachments. Dimension at most three is allowed for
every $n\geq3$ by E.1.

Kervaire's opened Theorem 1 \citep{Kervaire1969} is the related
manifold realization of finitely presented superperfect groups
in dimensions at least five. Our three-dimensional complex
argument is written out because that manifold theorem by itself
does not give the dimension required here.
\end{proof}

\paragraph{Outcome for Lead 4.}
Parts (b), (c), and (d) hold with the infinite version in E.5.
Part (a), as a claim about its two specified functors for every
simplicial set, is false in every dimension $n\geq2$; E.2 is the
correct replacement. The counterexample does not affect PRETOT
A.5, whose free-degeneracy hypothesis is present in its proof.
